% Options for packages loaded elsewhere
\PassOptionsToPackage{unicode}{hyperref}
\PassOptionsToPackage{hyphens}{url}
\documentclass[
  ignorenonframetext,
]{beamer}
\newif\ifbibliography
\usepackage{pgfpages}
\setbeamertemplate{caption}[numbered]
\setbeamertemplate{caption label separator}{: }
\setbeamercolor{caption name}{fg=normal text.fg}
\beamertemplatenavigationsymbolsempty
% remove section numbering
\setbeamertemplate{part page}{
  \centering
  \begin{beamercolorbox}[sep=16pt,center]{part title}
    \usebeamerfont{part title}\insertpart\par
  \end{beamercolorbox}
}
\setbeamertemplate{section page}{
  \centering
  \begin{beamercolorbox}[sep=12pt,center]{section title}
    \usebeamerfont{section title}\insertsection\par
  \end{beamercolorbox}
}
\setbeamertemplate{subsection page}{
  \centering
  \begin{beamercolorbox}[sep=8pt,center]{subsection title}
    \usebeamerfont{subsection title}\insertsubsection\par
  \end{beamercolorbox}
}
% Prevent slide breaks in the middle of a paragraph
\widowpenalties 1 10000
\raggedbottom
\AtBeginPart{
  \frame{\partpage}
}
\AtBeginSection{
  \ifbibliography
  \else
    \frame{\sectionpage}
  \fi
}
\AtBeginSubsection{
  \frame{\subsectionpage}
}
\usepackage{iftex}
\ifPDFTeX
  \usepackage[T1]{fontenc}
  \usepackage[utf8]{inputenc}
  \usepackage{textcomp} % provide euro and other symbols
\else % if luatex or xetex
  \usepackage{unicode-math} % this also loads fontspec
  \defaultfontfeatures{Scale=MatchLowercase}
  \defaultfontfeatures[\rmfamily]{Ligatures=TeX,Scale=1}
\fi
\usepackage{lmodern}
\ifPDFTeX\else
  % xetex/luatex font selection
\fi
% Use upquote if available, for straight quotes in verbatim environments
\IfFileExists{upquote.sty}{\usepackage{upquote}}{}
\IfFileExists{microtype.sty}{% use microtype if available
  \usepackage[]{microtype}
  \UseMicrotypeSet[protrusion]{basicmath} % disable protrusion for tt fonts
}{}
\makeatletter
\@ifundefined{KOMAClassName}{% if non-KOMA class
  \IfFileExists{parskip.sty}{%
    \usepackage{parskip}
  }{% else
    \setlength{\parindent}{0pt}
    \setlength{\parskip}{6pt plus 2pt minus 1pt}}
}{% if KOMA class
  \KOMAoptions{parskip=half}}
\makeatother
\setlength{\emergencystretch}{3em} % prevent overfull lines
\providecommand{\tightlist}{%
  \setlength{\itemsep}{0pt}\setlength{\parskip}{0pt}}
\usepackage{bookmark}
\IfFileExists{xurl.sty}{\usepackage{xurl}}{} % add URL line breaks if available
\urlstyle{same}
\hypersetup{
  hidelinks,
  pdfcreator={LaTeX via pandoc}}

\author{\texorpdfstring{}{}}
\date{}

\begin{document}

\begin{frame}[fragile]{Appendix D --- Inference Loop Mathematics}
\protect\phantomsection\label{appendix-d-inference-loop-mathematics}
This appendix formalizes the discrete-time active inference loop
executed by \texttt{cogant\ process} on the extracted A/B/C/D matrices
and reported in Section 5 (Table of VFE traces) and
\texttt{../cogant/docs/evaluation/EMPIRICAL\_CLAIM.md}. The formalism
follows Da Costa et al.~(2020) and the pymdp reference (Heins et al.,
2022); we restate it here in notation consistent with COGANT's
\texttt{cogant.process} module.

\begin{block}{D.1 POMDP formulation}
\protect\phantomsection\label{d.1-pomdp-formulation}
The extracted model is a discrete-time Partially Observable Markov
Decision Process
\texttt{(S,\ O,\ A,\ π,\ A\_mat,\ B\_mat,\ C\_mat,\ D\_mat)} with:

\begin{quote}
\textbf{S = \{s\_1, \ldots, s\_\textbar S\textbar\}} --- finite set of
hidden states. Cardinality is the product of factor cardinalities:
\texttt{\textbar{}S\textbar{}\ =\ ∏\_f\ \textbar{}S\_f\textbar{}}.

\textbf{O = \{o\_1, \ldots, o\_\textbar O\textbar\}} --- finite set of
observations. For multi-modality models,
\texttt{\textbar{}O\textbar{}\ =\ ∏\_m\ \textbar{}O\_m\textbar{}}.

\textbf{A ⊆ \{1, \ldots, \textbar A\textbar\}} --- finite set of
discrete actions (control states).

\textbf{π ∈ Π} --- policies, i.e.~finite sequences
\texttt{(a\_0,\ a\_1,\ …,\ a\_\{T−1\})\ ∈\ A\^{}T} over horizon
\texttt{T}.

\textbf{A\_mat ∈ ℝ\^{}\{\textbar O\textbar×\textbar S\textbar\}},
\texttt{A\_mat{[}o,\ s{]}\ =\ P(o\ \textbar{}\ s)} --- likelihood.

\textbf{B\_mat ∈
ℝ\^{}\{\textbar S\textbar×\textbar S\textbar×\textbar A\textbar\}},
\texttt{B\_mat{[}s\textquotesingle{},\ s,\ a{]}\ =\ P(s\textquotesingle{}\ \textbar{}\ s,\ a)}
--- state‑transition tensor.

\textbf{C\_mat ∈ ℝ\^{}\{\textbar O\textbar\}}, \texttt{C\_mat{[}o{]}}
--- log‑preference over observations.

\textbf{D\_mat ∈ ℝ\^{}\{\textbar S\textbar\}},
\texttt{D\_mat{[}s{]}\ =\ P(s\_0\ =\ s)} --- prior over initial states.
\end{quote}

All COGANT-extracted matrices satisfy the stochasticity conditions
\texttt{∑\_o\ A\_mat{[}o,\ s{]}\ =\ 1} for all \texttt{s} and
\texttt{∑\_\{s\textquotesingle{}\}\ B\_mat{[}s\textquotesingle{},\ s,\ a{]}\ =\ 1}
for all \texttt{(s,\ a)}; the GNN validator enforces these invariants at
emission time.
\end{block}

\begin{block}{D.2 Variational free energy functional}
\protect\phantomsection\label{d.2-variational-free-energy-functional}
Let \texttt{Q(s)} be an approximate posterior over hidden states and
\texttt{P(o,\ s)} the joint distribution defined by the generative model
\texttt{P(o,\ s)\ =\ A\_mat{[}o,\ s{]}\ ·\ D\_mat{[}s{]}}. The
variational free energy (VFE) is

\begin{quote}
\textbf{F{[}Q{]} = 𝔼\_\{Q(s)\}{[} log Q(s) − log P(o, s) {]}}

\begin{verbatim}
  = **KL( Q(s) ∥ P(s | o) ) − log P(o)**
\end{verbatim}
\end{quote}

The second equality (the ``Helmholtz decomposition'') shows that
minimizing \texttt{F} is equivalent to finding the posterior that best
approximates \texttt{P(s\ \textbar{}\ o)} up to a constant
\texttt{log\ P(o)} that depends only on the observation. Equivalently,

\begin{quote}
\textbf{F{[}Q{]} = 𝔼\_\{Q(s)\}{[}−log A\_mat{[}o, s{]}{]} − H{[}Q(s){]}
− 𝔼\_\{Q(s)\}{[}log D\_mat{[}s{]}{]}}
\end{quote}

decomposes VFE into three interpretable terms: the expected negative
log‑likelihood (prediction error), the negative entropy of the posterior
(ambiguity), and the expected log‑prior (complexity). COGANT's
\texttt{cogant\ process} computes this decomposition directly from the
extracted matrices.
\end{block}

\begin{block}{D.3 Variational inference via belief propagation}
\protect\phantomsection\label{d.3-variational-inference-via-belief-propagation}
For a single-factor discrete POMDP with observation \texttt{o\_t} at
time \texttt{t}, the posterior update is the normalized product

\begin{quote}
\textbf{Q(s\_t) ∝ A\_mat{[}o\_t, s\_t{]} · Q(s\_\{t\textbar t−1\})}
\end{quote}

where \texttt{Q(s\_\{t\textbar{}t−1\})} is the predicted state (the
result of applying the transition tensor to the previous posterior:
\texttt{Q(s\_\{t\textbar{}t−1\})\ =\ ∑\_\{s\_\{t−1\}\}\ B\_mat{[}s\_t,\ s\_\{t−1\},\ a\_\{t−1\}{]}\ ·\ Q(s\_\{t−1\})}).
Because the posterior is a categorical distribution over a finite set,
the update is exact --- there is no approximation --- and convergence of
the inner loop is trivial (single step). The belief propagation
terminology is retained because the formalism extends to factor-graph
inference when the hidden state is factorized into multiple independent
factors.
\end{block}

\begin{block}{D.4 VFE = 0.0 in the identity model}
\protect\phantomsection\label{d.4-vfe-0.0-in-the-identity-model}
The zoo/01\_simple\_state fixture demonstrates the identity case where
VFE converges to exactly zero. The extracted model has

\begin{quote}
\texttt{\textbar{}S\textbar{}\ =\ 1} (single factor, single cardinality
after aggregation)

\texttt{A\_mat\ =\ {[}{[}1.0{]}{]}} (identity likelihood)

\texttt{B\_mat{[}·,\ ·,\ a{]}\ =\ {[}{[}1.0{]}{]}} for all \texttt{a}
(identity transition, all actions)

\texttt{C\_mat\ =\ {[}0.0{]}} (no preference gradient)

\texttt{D\_mat\ =\ {[}1.0{]}} (fully certain prior)
\end{quote}

Substituting into the VFE decomposition:

\begin{quote}
\texttt{F\ =\ 𝔼\_\{Q(s)\}{[}−log\ A\_mat{[}o,\ s{]}{]}\ −\ H{[}Q(s){]}\ −\ 𝔼\_\{Q(s)\}{[}log\ D\_mat{[}s{]}{]}}

= \texttt{−log(1.0)\ −\ 0\ −\ log(1.0)}

= \textbf{0.0}
\end{quote}

The three terms vanish separately: the prediction error is zero because
\texttt{A\_mat{[}0,\ 0{]}\ =\ 1.0} and the observation is guaranteed,
the entropy is zero because \texttt{Q(s)\ =\ {[}1.0{]}} is a Dirac delta
on the single state, and the complexity term is zero because the prior
is also a Dirac. This is the correct and expected behaviour for any
fixture where the extracted model is a degenerate single-state POMDP;
the ten-step trace in
\texttt{../cogant/docs/evaluation/EMPIRICAL\_CLAIM.md} confirms
\texttt{F\ =\ −0.000000} at every step, which is the expected numerical
signature (the \texttt{−0} arises from the sign of
\texttt{log(1.0)\ =\ 0} after the negation in the prediction-error
term).
\end{block}

\begin{block}{D.5 Other regimes observed in the empirical runs}
\protect\phantomsection\label{d.5-other-regimes-observed-in-the-empirical-runs}
Three qualitatively distinct VFE regimes appear in the four zoo fixtures
reported in \texttt{../cogant/docs/evaluation/EMPIRICAL\_CLAIM.md}:

\begin{enumerate}
\item
  \textbf{\texttt{VFE\ =\ 0.0} (flat certainty).}
  \texttt{zoo/01\_simple\_state} and \texttt{zoo/02\_observer} ---
  identity A/B with \texttt{D\ =\ {[}1.0{]}}, no free energy gradient,
  no belief update happens because the prior is already exact.
\item
  \textbf{\texttt{VFE\ =\ 23.025851} (maximum uncertainty floor).}
  \texttt{zoo/04\_pomdp\_minimal} --- observation-only GNN where the
  likelihood matrix \texttt{A\_mat} is empty (the extracted model has no
  hidden-state factor). The runtime evaluates
  \texttt{−log(1e-10)\ =\ 23.025850929940457} as the floor for an
  unresolvable observation, which is the expected floor for the
  \texttt{cogant.process} implementation when the likelihood is
  vacuously defined.
\item
  \textbf{\texttt{VFE\ →\ 0.798508} (converging plateau).}
  \texttt{zoo/06\_hierarchical} --- two-factor hierarchical model with
  discriminative likelihood
  \texttt{A\_mat\ =\ {[}{[}0.9,\ 0.1{]},\ {[}0.1,\ 0.9{]}{]}}. The
  posterior collapses from the uniform prior
  \texttt{D\_mat\ =\ {[}0.5,\ 0.5{]}} to the certain state
  \texttt{Q(s)\ =\ {[}1.0,\ 0.0{]}} by \texttt{t\ =\ 2}; VFE rises from
  \texttt{F(t=0)\ =\ 0.751435} to the equilibrium
  \texttt{F(t≥4)\ =\ 0.798508}. The plateau value is the equilibrium
  free energy of the committed state under the \texttt{0.9\ /\ 0.1}
  likelihood and corresponds to the residual complexity term
  \texttt{−∑\_s\ Q(s)\ log\ D\_mat{[}s{]}} evaluated at the collapsed
  posterior.
\end{enumerate}
\end{block}

\begin{block}{D.6 Multi-episode D update rule and convergence}
\protect\phantomsection\label{d.6-multi-episode-d-update-rule-and-convergence}
For multi-episode runs, the prior \texttt{D\_mat} is updated via
empirical Bayes:

\begin{quote}
\textbf{D\_mat\^{}\{(k+1)\}{[}s{]} = α · D\_mat\^{}\{(k)\}{[}s{]} + (1 −
α) · 𝔼\_τ{[}Q\^{}\{(k)\}(s\_0){]}}
\end{quote}

where \texttt{α\ ∈\ {[}0,\ 1)} is a learning rate, \texttt{τ} indexes
episodes in the current batch, and \texttt{𝔼\_τ{[}Q\^{}\{(k)\}(s\_0){]}}
is the average initial posterior across episodes. The update is a convex
combination of the previous prior and the empirical distribution of
inferred initial states; since both sides lie on the probability simplex
and the mapping is a contraction (the average of a bounded distribution
is bounded), the iteration converges to a fixed point
\texttt{D\_mat\^{}*} at which
\texttt{D\_mat\^{}*\ =\ 𝔼\_τ{[}Q(s\_0\ \textbar{}\ D\_mat\^{}*){]}}.
Convergence rate is geometric with ratio \texttt{α}; in COGANT's default
configuration \texttt{α\ =\ 0.9}, so the D update takes on the order of
ten episodes to converge to within 10⁻³ of the fixed point. The update
is implemented in \texttt{cogant.process} as
\texttt{update\_prior\_from\_episodes(prior,\ episodes,\ alpha=0.9)} and
is disabled by default for the single-episode runs reported in
\texttt{../cogant/docs/evaluation/EMPIRICAL\_CLAIM.md}.
\end{block}

\begin{block}{D.7 Expected free energy and policy selection}
\protect\phantomsection\label{d.7-expected-free-energy-and-policy-selection}
For policy selection, COGANT uses the expected free energy (EFE) for
each candidate policy \texttt{π}:

\begin{quote}
**G(π) = ∑\_τ {[} 𝔼\_\{Q(o\_τ, s\_τ \textbar{} π)\}{[}log Q(s\_τ
\textbar{} π) − log P(o\_τ, s\_τ){]} {]}**

\begin{verbatim}
   = **∑_τ [ risk(π, τ) + ambiguity(π, τ) ]**
\end{verbatim}
\end{quote}

where \texttt{risk} is the KL divergence between predicted observations
and preferences (\texttt{C\_mat}) and \texttt{ambiguity} is the expected
entropy of the likelihood under predicted states. The implementation in
\texttt{cogant.process.evaluate\_policies} computes \texttt{G(π)} for
every policy in the finite policy space and selects the argmin (softmax
with temperature = 0 in the deterministic default). On
zoo/01\_simple\_state with \texttt{C\_mat\ =\ {[}0.0{]}}, both
\texttt{u\_c0} and \texttt{u\_c1} score \texttt{G\ =\ 0.0} identically;
the argmin tie-break returns \texttt{u\_c0} every step, which is the
behaviour observed in
\texttt{../cogant/docs/evaluation/EMPIRICAL\_CLAIM.md} §3.
\end{block}
\end{frame}

\begin{frame}
\end{frame}

\end{document}
